%%%PAGE 4 rot=0 legibility=good summary=衰变+动量(α衰变两圆外切/β衰变两圆内切及相关公式)、光谱分类思维导图、巴耳末公式与里德伯常量
%%%BLOCK id=004-01 ch=17 sec=衰变+动量·磁场中径迹 kind=method alt=07,11 cont=none y=0.02-0.42
\textbf{$\langle$衰变+动量$\rangle$} % 原稿：标题被一对大尖括号圈起

\noindent $\alpha$衰变.\quad $\nuc{A}{Z}{X}\to\nuc{A-4}{Z-2}{Y}+\nuc{4}{2}{He}$\quad \fbox{退2\unclear{少}4} % 圈内文字“退2少4”中“少”字不清(也可能是“以”)

两圆外切，$\alpha$粒子半径较大

%FIG p004_a 0.268 0.14 0.375 0.20
\notefig[0.25]{figures/p004_a.png}
\begin{gather*}
m_1v_1-m_2v_2=0\\
E=\frac{p^2}{2m_1}+\frac{p^2}{2m_2}\\
r=\frac{p}{qB}
\end{gather*}
\[
I=\frac{q}{T}=\frac{q\,qB}{2\pi m}=\frac{q^2B}{2\pi m}
\]

\noindent $\beta$衰变\quad $\nuc{A}{Z}{X}\to\nuc{A}{Z+1}{Y}+\nuc{0}{-1}{e}$\quad \fbox{进1\unclear{位}} % 圈内文字“进1位”中“位”字不清

两圆内切，$\beta$粒子半径较大.

%FIG p004_b 0.275 0.331 0.395 0.41
\notefig[0.25]{figures/p004_b.png}
%%%END
%%%BLOCK id=004-02 ch=17 sec=光谱·分类 kind=summary alt=none cont=none y=0.50-0.69
\textbf{光谱}
\begin{itemize}
\item 发射光谱
  \begin{itemize}
  \item 连续谱
  \item \fbox{线状谱}
  \end{itemize}
\item \fbox{吸收光谱}
\end{itemize}
\noindent \fbox{线状谱}、\fbox{吸收光谱}$\;\to\;$原子的特征谱线 % 原稿：“线状谱”与“吸收光谱”的引线合并后，以一个箭头指向“原子的特征谱线”
%%%END
%%%BLOCK id=004-03 ch=17 sec=氢原子光谱·巴耳末公式 kind=formula alt=none cont=none y=0.67-0.73
\[
\frac{1}{\lambda}=R\left(\frac{1}{2^2}-\frac{1}{n^2}\right)\qquad\bigl(n=3,4,5\ldots\quad \overset{\text{里德伯常量}}{R=1.10\times10^{7}\unit{m^{-1}}}\bigr)
\] % CHECK: 第二项分母手写形似“2^n”，按巴耳末公式转为 n^2；“里德伯常量”写在 R=… 上方
%%%END
%%%PAGEEND 4
